\documentclass[11pt]{article}

% ---- Packages (all standard in a full TeX Live / MiKTeX install) ----
\usepackage[utf8]{inputenc}
\usepackage[T1]{fontenc}
\usepackage{lmodern}
\usepackage[margin=1in]{geometry}
\usepackage{amsmath,amssymb,amsthm}
\usepackage{mathtools}
\usepackage{booktabs}
\usepackage{array}
\usepackage{multirow}
\usepackage{microtype}
\usepackage{graphicx}
\usepackage{xcolor}
\usepackage{listings}
\usepackage[hidelinks]{hyperref}
\usepackage[capitalize,noabbrev]{cleveref}

% ---- Theorem environments ----
\theoremstyle{plain}
\newtheorem{theorem}{Theorem}
\newtheorem{lemma}{Lemma}
\newtheorem{proposition}{Proposition}
\theoremstyle{definition}
\newtheorem{definition}{Definition}
\newtheorem{construction}{Construction}
\theoremstyle{remark}
\newtheorem{remark}{Remark}

% ---- Macros ----
\newcommand{\Enc}{\mathsf{Enc}}
\newcommand{\Dec}{\mathsf{Dec}}
\newcommand{\Fabb}{\mathcal{F}_{\mathsf{ABB}}}
\newcommand{\ReLU}{\mathrm{ReLU}}
\newcommand{\Z}{\mathbb{Z}}
\newcommand{\bt}[1]{\texttt{#1}}
\newcommand{\toMPC}{\mathsf{fhe{\to}shares}}
\newcommand{\toFHE}{\mathsf{shares{\to}fhe}}
\DeclareMathOperator{\poly}{poly}
\DeclareMathOperator{\maj}{maj}

\lstset{basicstyle=\ttfamily\small,columns=fullflexible,keepspaces=true,
        breaklines=true,frame=none,xleftmargin=1em}

\title{\bfseries The Switch Is the Bootstrap:\\
Deep Private Inference at Zero FHE Bootstraps\\
via Scheme Switching}

\author{%
  (author names elided for review)%
}
\date{}

\begin{document}
\maketitle

\begin{abstract}
Hybrid private-inference systems evaluate a neural network's linear layers
under homomorphic encryption (FHE) and its nonlinearities under secure
multi-party computation (MPC), switching representation between the two.
Prior work treats this scheme switch as an engineering necessity. We observe
that it is also a \emph{noise-reset mechanism}: the re-encryption that
returns an activation from MPC shares to an FHE ciphertext produces a
ciphertext distributed identically to a \emph{fresh}, depth-one encryption.
Consequently a network of \emph{arbitrary depth} is evaluated with
\textbf{zero FHE bootstraps} and cost \textbf{linear in depth}, using an FHE
engine that supports only a \emph{single} homomorphic multiplication and
\emph{no evaluation keys}. We make the argument precise (\Cref{thm:main}),
instantiate the FHE side so that a full matrix--vector product with
\emph{both} the weights and the activations encrypted is one
ciphertext$\times$ciphertext multiplication decrypted under $s^2$ (no
relinearization), and give a convolution packing that yields an entire
output channel per multiplication. The MPC side is a triple-free,
communication-batched honest-majority protocol with constant-round garbled
nonlinearities. We implement the whole pipeline and validate it by
\emph{bit-exact} agreement with a plaintext integer reference at every depth.
On real biomedical images (MedMNIST v2), a fully encrypted two-convolution
CNN --- four encrypted linear layers, three MPC nonlinearities, three scheme
switches, and zero bootstraps --- diagnoses a pediatric chest X-ray in
$\mathbf{2.9}$ seconds at $84.78\%$ accuracy, and reaches $98.89\%$ on MNIST,
with every encrypted prediction provably equal to its plaintext counterpart.
Unlike two-party inference services, the model weights are encrypted as well
as the input. Getting from $32.7$ s to $2.9$ s cost us a diagnosis of our
own, which we report: the dominant cost was neither the FHE arithmetic nor
the MPC protocols but the \emph{plumbing between them} --- an unbatched
scheme-switch import ($8.3\times$), and then the query-independent masking
that same switch performs ($1.7\times$, moved offline). The construction
never changed; only its round structure did.
\end{abstract}

\section{Introduction}\label{sec:intro}

Fully homomorphic encryption~\cite{Gentry} lets a server compute on
encrypted data, but every homomorphic multiplication grows ciphertext noise;
after a bounded number of multiplications the ciphertext must be
\emph{bootstrapped} --- an expensive homomorphic decryption --- to continue.
Deep networks therefore incur bootstrapping cost that scales with depth.

An orthogonal line of work evaluates the \emph{linear} layers under FHE and
the \emph{nonlinear} layers (ReLU, max-pooling, argmax) under
MPC~\cite{Gazelle,Delphi,Cheetah,CrypTFlow2}, because comparisons are cheap
in MPC and expensive in FHE. Between the two representations sits a
\emph{scheme switch}: mask the ciphertext, decrypt the mask, compute the
nonlinearity on additive secret shares, and re-encrypt the result. These
systems present the switch as unavoidable overhead.

\paragraph{Our observation.} The return leg of the switch --- re-encrypting
each party's share component and summing --- produces a ciphertext that is a
sum of \emph{fresh} encryptions. A sum of fresh encryptions is itself a
fresh encryption: it carries only fresh noise and sits at multiplicative
depth one, indistinguishable from a first-layer input. The switch that a
system already performs for the ReLU therefore \emph{also refreshes the FHE
ciphertext}, and it does so for free. The per-layer FHE noise budget is a
constant \emph{independent of network depth}, so bootstrapping is never
needed --- not deferred, but absent. We call this ``the switch is the
bootstrap.''

This inverts a familiar tradeoff. Because the FHE engine only ever performs
one multiplication before the next switch resets it, we may use an FHE
scheme that is fast precisely \emph{because} it supports only depth one and
requires \emph{no evaluation keys} (no relinearization keys, no Galois keys,
no bootstrapping keys). A ciphertext$\times$ciphertext product is a degree-2
ciphertext; rather than relinearizing it back to degree one, we simply
decrypt it under $s^2$ at the next switch --- which opens a masked value
anyway. The result is a private-inference construction whose FHE cost per
layer is a single multiplication and whose total FHE cost is linear in depth
with a small constant.

\paragraph{Contributions.}
\begin{enumerate}
\item \textbf{Switch-as-bootstrap (\Cref{sec:construction}).} We formalize
the claim that the MPC round trip refreshes the FHE ciphertext to a
depth-one state, giving a per-layer noise bound independent of depth and
hence unbounded-depth evaluation at zero bootstraps. The argument is
scheme-agnostic: any additively-homomorphic FHE with a cheap depth-one
multiply instantiates it.

\item \textbf{Evaluation-key-free, both-operands-private linear layers
(\Cref{sec:fhe}).} Instantiating the FHE side with RLWE (BFV) and
reversed-coefficient packing, a full matrix--vector product with both the
weight matrix \emph{and} the activation vector encrypted is a single
ciphertext$\times$ciphertext multiplication decrypted under $s^2$, with no
relinearization and no evaluation keys. A Cheetah-style convolution packing
yields an entire output channel per multiplication. This strengthens the
privacy model of prior hybrid systems, which keep the model in plaintext.

\item \textbf{A triple-free, batched honest-majority MPC layer
(\Cref{sec:mpc}).} Multiplication uses ABY3-style resharing with PRF-derived
correlated randomness --- one ring element per party per multiplication, no
Beaver triples, no dealer. Bit decomposition is a Kogge--Stone construction
batched across an entire layer, so a ReLU or max-pooling layer costs a
constant number of communication rounds regardless of width.

\item \textbf{Composability (\Cref{sec:uc}).} The pipeline is an Arithmetic
Black Box; a small gadget library (\bt{relu}, \bt{max2}, \bt{maxpool2d},
\bt{argmax}, \bt{dense}, \bt{conv2d}) realizes standard network layers as
pure functions of the ABB primitives, so any architecture built from them
inherits security and correctness with no per-architecture proof.

\item \textbf{An implementation validated bit-exactly
(\Cref{sec:eval}).} We implement the entire construction and check it against
a plaintext integer reference: the encrypted output equals the plaintext
output \emph{bit for bit} at every measured depth. On real MNIST test
digits, a fully encrypted two-convolution CNN reaches $98.89\%$ accuracy at
zero bootstraps.

\item \textbf{Private diagnosis on real medical images
(\Cref{sec:diagnosis}), and a corrected cost model.} The same pipeline,
unchanged, classifies MedMNIST v2~\cite{MedMNIST} chest X-rays and abdominal
CT at $84.78\%$ / $81.18\%$ --- within $0.2$ points of the float network it
was quantized from --- in $4.9$ s per encrypted diagnosis. Reaching that
number required finding that the pipeline's dominant cost was not what a
segment-level measurement suggested: $88\%$ of end-to-end runtime was an
unbatched scheme-switch \emph{import}, not the FHE arithmetic and not the MPC
comparison protocols. Batching it ($8.3\times$), garbling the
nonlinearities ($1.7\times$ on the MPC segment), and moving the switch's
query-independent masking offline ($1.7\times$ again, \Cref{sec:preproc})
give $2.9$ s, of which $73\%$ is now the FHE multiplications themselves.
\end{enumerate}

\paragraph{Honesty about scope.} We separate what is proven, what is
implemented, and what is simulated. The design is semi-honest,
honest-majority. Two aspects of the \emph{implementation} stand in for their
deployable form and are disclosed as limitations (\Cref{sec:security}): the
masked decryption uses a single client key as a placeholder for threshold
FHE, and the fast reference FHE engine (scale-invariant NTRU) is
\emph{insecure} at the parameters this construction needs. The secure
instantiation is RLWE, on which all headline accuracy and correctness
numbers are measured; NTRU is retained only as a speed reference for
isolating packing cost from FHE-scheme cost.

\section{Preliminaries}\label{sec:prelim}

\paragraph{Notation.} We work over the polynomial ring
$R_q = \Z_q[X]/(X^n+1)$ with $n$ a power of two. Lowercase letters denote
ring elements; $s$ is an RLWE secret key. For an integer $v$ and width $b$,
$\langle v\rangle$ denotes an additive secret sharing over $\Z_{2^b}$. We
write $x \gg s$ for a logical right shift by $s$ bits.

\paragraph{RLWE encryption (BFV).}\label{par:bfv}
We use the Brakerski/Fan--Vercauteren scheme~\cite{FV}. A plaintext is a
polynomial with coefficients mod $t$; a ciphertext is a pair
$(c_0,c_1)\in R_q^2$ with $c_0 + c_1 s \approx \Delta m$ for
$\Delta=\lfloor q/t\rfloor$. Addition is componentwise. The product of two
ciphertexts is a \emph{degree-2} ciphertext $(d_0,d_1,d_2)$ satisfying
$d_0 + d_1 s + d_2 s^2 \approx \Delta^2 m_1 m_2 / \Delta$ after rescaling;
usually one \emph{relinearizes} $(d_0,d_1,d_2)$ back to a pair using an
evaluation key. We never relinearize: the next scheme switch decrypts the
degree-2 ciphertext directly under $s^2$
(\Cref{sec:fhe}). We take $n=4096$ and $\log_2 q = 100$, at or below the
$128$-bit security ceiling of $\log_2 q \le 109$ for
$n=4096$~\cite{HEStandard}.

\paragraph{Replicated secret sharing (ABY3).}\label{par:rss}
For the honest-majority three-party segment we use $2$-out-of-$3$ replicated
sharing~\cite{ABY3}: a value $v=v_0+v_1+v_2$ over $\Z_{2^b}$ is shared so
party $i$ holds $(v_i, v_{i+1})$. Addition and multiplication-by-constant
are local. Multiplication is described in \Cref{sec:mpc}.

\paragraph{Scheme switching.} A \emph{switch to MPC}, $\toMPC$, takes an FHE
ciphertext of a value $x$, has the parties add fresh random masks, decrypts
the single masked ciphertext, and produces an additive sharing
$\langle x\rangle$; only a uniformly-masked value is ever opened. A
\emph{switch to FHE}, $\toFHE$, has each party freshly encrypt its share
component; the evaluator sums the ciphertexts.

\section{The Switch-as-Bootstrap Construction}\label{sec:construction}

\subsection{The layer template}

A network is a sequence of layers $x \mapsto \sigma(L(x))$, where $L$ is an
affine map (dense or convolution, possibly with encrypted weights) and
$\sigma$ is a nonlinearity (ReLU, optionally followed by max-pooling). Each
layer is evaluated as:
\begin{enumerate}
\item \textbf{FHE linear map.} Given a fresh ciphertext $c=\Enc(x)$, compute
$c' = L(c)$: one homomorphic multiplication (\Cref{sec:fhe}). $c'$ is a
degree-2 ciphertext with depth-one noise.
\item \textbf{Switch to MPC.} $\langle y\rangle \leftarrow \toMPC(c')$,
opening only a masked value. The ciphertext noise is now discarded.
\item \textbf{Nonlinearity in MPC.}
$\langle a\rangle \leftarrow \sigma(\langle y\rangle)$ (\Cref{sec:mpc}).
\item \textbf{Switch to FHE.} $c_{\mathrm{next}} \leftarrow \toFHE(\langle
a\rangle)$: a fresh ciphertext for the next layer.
\end{enumerate}

\subsection{Main claim}

\begin{theorem}[Switch-as-bootstrap]\label{thm:main}
Let $\Enc$ be an additively homomorphic encryption scheme whose fresh
ciphertexts carry noise at most $\nu_0$ and admit one homomorphic
multiplication within a decryption-correct noise budget $B$ (i.e.
$\poly(\nu_0)\le B$). In the layer template above, the ciphertext entering
the FHE step of \emph{every} layer is distributed identically to a fresh
encryption: it has multiplicative depth one and noise at most $w\cdot\nu_0$,
where $w$ is the layer's fan-in width. Consequently, for a network of any
depth $D$, decryption correctness holds at every layer under the single
budget $B$, and no bootstrapping is invoked.
\end{theorem}

\begin{proof}[Proof sketch]
The claim is an invariant maintained by induction on the layer index. The
base case is the network input, a fresh encryption by construction. For the
inductive step, assume layer $\ell$ begins with a fresh ciphertext. Steps
1--3 consume it and produce an additive sharing $\langle a\rangle$ living
entirely in the MPC domain; the FHE ciphertext and its noise are discarded
at Step 2. Step 4 forms
$c_{\mathrm{next}} = \sum_{i} \Enc(a^{(i)})$, a sum of freshly sampled
encryptions of the parties' share components $a^{(i)}$. By additive
homomorphism $c_{\mathrm{next}}$ encrypts $\sum_i a^{(i)} = a$, and because
each summand is an independent fresh encryption, $c_{\mathrm{next}}$ is
distributed as a fresh encryption of $a$ with noise bounded by the sum of at
most $w$ fresh noises, i.e. $w\cdot\nu_0$, at multiplicative depth one. Thus
layer $\ell+1$ begins in the same fresh state. Since the state is fresh at
every layer, the single-multiplication budget $B$ suffices at every layer,
independent of $D$. Bootstrapping, the operation used only to reset
accumulated noise, is therefore never invoked.
\end{proof}

\begin{remark}[Cost linear in depth]
Each layer performs a constant amount of FHE work (one multiply, $O(1)$ mask
ciphertexts), one masked decryption, and one MPC nonlinearity. Total cost is
therefore linear in depth, with no super-linear bootstrapping term. This is
the qualitative separation from leveled FHE, where the per-multiply cost
grows with the remaining level budget or a bootstrap is amortized in.
\end{remark}

\begin{remark}[Why we can drop evaluation keys]
\Cref{thm:main} needs only \emph{one} multiplication per layer and never
composes two multiplications homomorphically. Relinearization exists to make
a degree-2 product usable by a \emph{subsequent} homomorphic operation; here
the subsequent operation is a decryption (the switch), which handles a
degree-2 ciphertext natively by evaluating $d_0+d_1 s + d_2 s^2$. No
relinearization key is needed. Coefficient packing (\Cref{sec:fhe})
avoids slot rotations, so no Galois keys are needed either. And no
bootstrapping key is needed by the theorem. The FHE public material is a
single public key.
\end{remark}

\subsection{Correctness of the fixed-point values}

\Cref{thm:main} concerns ciphertext \emph{noise}. Independently, the integer
activations must not overflow the share modulus $2^b$. We use a fused
nonlinearity
\begin{equation}\label{eq:fused}
a' \;=\; \ReLU\!\big((W\!\cdot\! a + (\beta \ll s)) \gg s\big),
\end{equation}
where the arithmetic right shift by $s$ rescales the fixed-point product so
that activation magnitude stays roughly constant across layers. Two points
make this exact.

First, the rescale is a \emph{logical} shift, yet $W\!\cdot\!a+(\beta\ll s)$
can be negative before the ReLU. Correctness holds because the shifted value
is retained only on the non-negative branch of the ReLU --- where a logical
and an arithmetic shift agree --- and zeroed on the negative branch. Hence a
single bit decomposition of the pre-activation yields \emph{both} the sign
bit (for the ReLU) and the shifted magnitude (for the rescale): one bit
decomposition per layer, not two.

Second, the construction is \emph{exact} integer fixed-point. The encrypted
output equals the plaintext integer network bit for bit --- not up to an
approximation error --- which is what our experiments verify as
$\texttt{correct}=\texttt{true}$ (\Cref{sec:eval}). This is a much stronger
correctness statement than the approximate equality reported by
CKKS-based approximate inference.

\section{Evaluation-Key-Free Linear Layers}\label{sec:fhe}

\subsection{Both operands encrypted, no relinearization}

We pack a length-$m$ vector into the coefficients of a degree-$m$ polynomial.
To compute an inner product $\langle u, v\rangle$, we pack $u$ in natural
order and $v$ in \emph{reversed} order; the degree-$(m-1)$ coefficient of the
polynomial product $u(X)\cdot v(X)$ is exactly $\sum_i u_i v_i$. Both
polynomials are ciphertexts, so a matrix--vector product is a
ciphertext$\times$ciphertext multiplication in which \emph{both} the weights
and the activations are encrypted. The product is a degree-2 ciphertext,
which the next switch decrypts under $s^2$; we never relinearize.

For a matrix with $\rho$ output rows and inner dimension $m$, the reversed
layout places row $r$'s weights at coefficient offset $r\cdot m$, so a single
degree-$n$ product yields up to $\lfloor n/m\rfloor - 1$ separable outputs; a
layer of $\rho$ rows costs $\lceil \rho m/n\rceil$ multiplications. We show
this is optimal for the coefficient encoding: a degree-$n$ product admits at
most $n/m$ separable outputs because distinct rows require disjoint weight
windows. SIMD-slot packing matches the throughput but requires rotations and
Galois keys, forfeiting the evaluation-key-free property.

\subsection{Exact negacyclic products via a multi-modular NTT}

For a power-of-two modulus $q$ the ring $R_q$ has no primitive $2n$-th root
of unity, so no direct NTT exists. We instead compute the \emph{exact}
integer negacyclic product using several NTT-friendly primes as a
computational device and reconstruct the integer coefficients by the Chinese
Remainder Theorem (Garner's algorithm), then reduce mod $q$ (and apply the
rounded-shift rescale). This is differentially tested against a schoolbook
reference at full coefficient width and $n=4096$. It accelerates the BFV
multiplication by roughly $11\times$ over the schoolbook implementation it
replaces.

\subsection{Convolution packing: an output channel per multiplication}

Polynomial multiplication is itself a convolution, which we exploit
directly. Flatten a $C\times H\times W$ input row-major into coefficients and
pack a kernel \emph{reversed on every axis} (channel, row, column) at the
input's strides. Then a single ciphertext$\times$ciphertext product places
the valid convolution output for an entire output channel at a fixed set of
coefficients. Two facts make it exact:
\begin{itemize}
\item \textbf{Channels pair off.} The reversed channel offset makes
$c + (C{-}1{-}c') = C{-}1$ hold only when $c=c'$, so every input channel's
contribution lands on the same output coefficient and they sum --- exactly
what a multi-channel convolution requires.
\item \textbf{No row-wrap contamination.} A neighbouring row could only leak
into an output coefficient via a column carry, which would require a
negative output column index --- impossible for a valid output. The read-out
coefficients are therefore clean.
\end{itemize}
The layer costs $\mathrm{out\_c}$ multiplications, \emph{independent of image
size}, versus $\lceil \mathrm{out\_c}\cdot o_h o_w\cdot CHW / n\rceil$ for an
im2col/Toeplitz formulation. At $28\times28$ with a $3\times3$ kernel and two
output channels this is $2$ multiplications instead of $1352$.

\section{The MPC Layer}\label{sec:mpc}

\subsection{Triple-free multiplication}

With replicated sharing $x=x_0+x_1+x_2$, each party $i$ holds $(x_i,
x_{i+1})$ and can form three of the nine cross terms of $x\cdot y$ locally:
\begin{equation}\label{eq:reshare}
z_i \;=\; x_i y_i + x_i y_{i+1} + x_{i+1} y_i + \alpha_i,
\end{equation}
where the $\alpha_i$ are zero-sum correlated randomness derived locally from
pairwise PRF keys, $\alpha_i = F(k_i, \mathrm{ctr}) - F(k_{i-1},
\mathrm{ctr})$. Summing over the three parties covers all nine cross terms
exactly once, so $\sum_i z_i = (\sum_i x_i)(\sum_i y_i) = xy$; the $\alpha_i$
cancel in the sum but re-randomize each party's share so nothing leaks. Each
multiplication is one ring element sent to one neighbour --- no Beaver
triples~\cite{Beaver}, no dealer, nothing precomputed beyond a one-time key
handshake. The PRF counter is a globally distinct value per multiplication,
so no two products ever share an $\alpha$ (which would expose their
difference).

\subsection{Batched bit decomposition}

Comparisons (for ReLU, max-pooling, argmax) reduce to extracting the sign
bit of a shared value, which requires an arithmetic-to-binary decomposition.
We use the ABY3 decomposition: (i) binary-share the three additive
components for free (each is already known to exactly the two parties a
replicated sharing needs); (ii) reduce three shares to two with a carry-save
adder, exploiting that in this encoding $x\oplus y\oplus z$ and
$\maj(x,y,z)$ fall out of the \emph{same} four products (two rounds, not
four); (iii) add the remaining two with a Kogge--Stone parallel-prefix
adder~\cite{KoggeStone}, whose $(g,p)$ combine
$G = g \oplus (p \wedge g')$ is local because a group cannot both generate
and propagate a carry. The depth is $2 + 1 + \lceil\log_2 b\rceil + 1
\approx 9$ rounds at $b=20$.

Crucially, we \emph{batch} the decomposition across an entire layer: every
internal multiplication round is flattened over all lanes, so decomposing a
whole ReLU or max-pooling layer costs the same $\approx 9$ rounds as
decomposing a single value. Without batching a $1352$-neuron ReLU layer
would run $\approx 12{,}000$ sequential rounds; with it, $\approx 9$. This
single change is a $6.7\times$ end-to-end speedup on the MPC segment
(\Cref{sec:eval}).

\section{Composability}\label{sec:uc}

The construction manipulates secret values only through opaque handles
(a share reference or a ciphertext handle) and a fixed set of primitive
operations; it never inspects the representation. This is exactly the
Arithmetic Black Box abstraction~\cite{DamgardNielsen,SPDZ}: an ideal
functionality $\Fabb$ storing values by handle and exposing
$\{\mathsf{add}, \mathsf{mul}, \mathsf{lt}, \mathsf{share\_public},
\mathsf{reveal}\}$ plus an internal, invisible representation switch. The
scheme switch of \Cref{sec:construction} is that internal operation; it is
never observable at the ABB interface.

A small gadget library realizes standard layers as pure functions of the
ABB primitives (\Cref{tab:gadgets}). Because each gadget is a composition of
ABB operations, if each primitive UC-realizes~\cite{Canetti} its $\Fabb$
operation, each gadget UC-realizes the corresponding ideal gadget, and any
network --- any sequence of gadgets --- UC-realizes the ideal inference
functionality; the simulator for a network is the composition of the
per-gadget simulators. We therefore do not prove ``private CNN inference''
as a monolith: we exhibit a handful of gadgets and compose them. Swapping in
more channels, more convolution blocks, or residual adds requires no new
security argument, only a new sequence of gadget calls.

\begin{table}[t]
\centering
\caption{The gadget library as compositions of ABB primitives.}
\label{tab:gadgets}
\begin{tabular}{@{}lll@{}}
\toprule
gadget & definition & primitives \\
\midrule
\bt{relu} & $x\cdot(0<x)$ & \bt{sub}, \bt{lt}, \bt{mul} \\
\bt{max2} & $\mathsf{mux}(a<b,\,b,\,a)$ & \bt{lt}, \bt{sub}, \bt{mul}, \bt{add} \\
\bt{maxpool2d} & fold of \bt{max2} over a window & (above) \\
\bt{argmax} & scan tracking $(\text{best},\text{idx})$ & \bt{lt}, \bt{mul}, \bt{add}, \bt{share\_public} \\
\bt{dense} / \bt{conv2d} & $\sum W\!\cdot\! x + b$ & \bt{mul}, \bt{add} \\
\bottomrule
\end{tabular}
\end{table}

\section{Security Model}\label{sec:security}

\paragraph{Threat model.} Semi-honest adversary; honest-majority
three-party MPC for the nonlinear segment; the FHE evaluator is an untrusted
server. Only the final output (the class index) is revealed. Intermediate
activations remain shared or encrypted, and each switch opens only a
uniformly-masked value, which is simulatable.

\paragraph{Disclosed implementation gaps.} We distinguish the design from
its current implementation.
\begin{enumerate}
\item \textbf{Threshold decryption is simulated.} The masked decryption uses
a single client key as a stand-in for threshold FHE~\cite{AsharovMHE}. The
data flow is already correct --- only uniformly-masked values are opened ---
so a deployable instantiation secret-shares the FHE key without changing the
protocol logic.
\item \textbf{The fast reference FHE engine is insecure; the secure one is
RLWE.} We port scale-invariant NTRU~\cite{LTV,YASHE} as a speed reference
only. It is overstretched even at depth one: the NTRU fatigue
point~\cite{ABDucas,KirchnerFouque,DucasvanWoerden} for the parameters a
usable plaintext modulus forces is well below the required $q$, so we do not
claim it secure. The secure instantiation is RLWE (BFV), which has no NTRU
lattice and no fatigue point, and supports the identical depth-one,
no-relinearization trick. \emph{All headline accuracy and correctness numbers
are measured on BFV.} A machine-checked regression asserts the ring
dimension and modulus are within the standardized $128$-bit
budget~\cite{HEStandard}.
\item \textbf{Malicious security is out of scope}, as is standard for this
line of work.
\end{enumerate}

\section{Evaluation}\label{sec:eval}

We implement the full construction in a single Rust crate (no GPU). The
three MPC parties run a real replicated-sharing wire protocol over serialized
messages, with a network-cost model (round-trip time and bandwidth). Every
benchmark verifies the encrypted result against a plaintext integer reference
and reports \texttt{correct}; all numbers are emitted by the tool, not
hand-entered. Wall-clock figures are single-thread on commodity hardware and
are meant to be read as ratios, not absolutes.

\subsection{Correctness and the zero-bootstrap invariant}

At every measured depth and width, the encrypted class equals the plaintext
class \emph{bit for bit}, and the bootstrap counter is exactly zero. This is
the operational content of \Cref{thm:main}: had noise accumulated with depth,
deep outputs would diverge from the plaintext reference; they do not. The
same invariant holds for the fully encrypted two-convolution CNN of
\Cref{sec:mnist}, whose four encrypted linear layers and three switches all
report zero bootstraps.

\subsection{The MPC segment: making it private, then fast}

\Cref{tab:mpc} traces the MPC segment of the flagship CNN ($28\times28$
input) through the sequence of protocol changes described in
\Cref{sec:mpc}. An initial masked-open comparison was found to leak its input
(the mask carried no secret entropy) and was replaced by the ABY3 bit
decomposition; correctness held throughout, but the correct protocol costs
more. Eliminating Beaver triples (PRF resharing) then cut bandwidth
$15\times$, and batching the decomposition across each layer cut time a
further $6.7\times$.

The last row is the one worth dwelling on, because it says the previous four
were optimizing the wrong thing. Having batched every \emph{comparison}, we
instrumented the segment per phase and found that the comparisons were
never its cost: $19.1$ of the remaining $20.4$ s was the \emph{import} ---
the MPC-side half of the scheme switch, which reconstructs an opened layer
into secret shares. Importing ran three blocking input-sharing rounds
\emph{per activation}, so the $5408$-wide first convolution cost $16{,}224$
sequential round trips, while the ReLU those values were being imported for
cost $0.6$ s. Batching the import per layer --- the same protocol, the same
randomness, the same bytes, in three rounds instead of $3n$ --- takes it to
$8$ ms and the end-to-end pipeline from $21.8$ s to $2.6$ s.

\begin{table}[t]
\centering
\caption{MPC segment of the $28\times28$ CNN across protocol revisions. Every
row is bit-exact against the plaintext reference; only privacy and cost
differ. The last two rows change no protocol --- only how many round trips
carry the same bytes. Measured at zero latency, which is where the tournament
\bt{argmax} looks least impressive: its win is in rounds, and at RTT $20$ ms
it is $2.4\times$ on that phase (\Cref{sec:gc}).}
\label{tab:mpc}
\begin{tabular}{@{}lrr@{}}
\toprule
protocol stage & MPC time & wire bytes \\
\midrule
leaky masked-open (insecure baseline) & $17.6$ s & $1.2$ MB \\
correct bit decomposition, Beaver triples & $76.9$ s & $159.4$ MB \\
\quad + triple-free PRF resharing & $40.8$ s & $10.4$ MB \\
\quad + batched ReLU decomposition & $20.7$ s & $10.2$ MB \\
\quad + batched max-pooling & $6.1$ s & $10.0$ MB \\
\quad + batched scheme-switch import & $1.2$ s & $10.0$ MB \\
\quad + tournament \bt{argmax} & $\mathbf{1.15}$ s & $10.0$ MB \\
\bottomrule
\end{tabular}
\end{table}

We report this misdiagnosis rather than quietly fixing it, because the shape
of the error is a general hazard for hybrid systems. ``The MPC segment
dominates'' was true; ``therefore the MPC \emph{protocols} are the
bottleneck'' did not follow. The segment's cost lived in the plumbing between
the two representations --- precisely the part that a construction organized
around scheme switching performs most often, and the part that neither the
FHE nor the MPC literature owns. A per-phase breakdown, not a per-segment
one, is what makes such a cost visible.

\subsection{Comparisons: garbled circuits vs.\ bit decomposition}\label{sec:gc}

Because comparisons dominate the MPC segment, we ask whether garbled circuits
--- constant-round --- beat the log-round bit decomposition. We built a real
half-gates garbling engine (free-XOR~\cite{FreeXOR} offset with lsb one and
Zahur--Rosulek--Evans half-gate AND gates~\cite{HalfGates}, two ciphertexts
per gate) and a $\mathrm{MSB}(x+y)$ comparator, bit-exact against the
plaintext comparison, and compared it head-to-head with the batched bit
decomposition. Crucially, \emph{both} run through the identical transport ---
real serialization, message passing, latency injection and byte counting ---
so there is no measurement asymmetry.

\begin{table}[t]
\centering
\caption{One layer of $1352$ comparisons at $20$ bits: garbled circuit vs.\
batched bit decomposition, measured through the identical transport. Both
bit-exact.}
\label{tab:gc}
\begin{tabular}{@{}lrrr@{}}
\toprule
network & a2b & GC & GC speedup \\
\midrule
structural & 9 rd, $5.9$ MB & 1 rd, $1.7$ MB & --- \\
LAN (RTT $0.1$ ms, $10$ Gbps) & $279$ ms & $66$ ms & $4.2\times$ \\
metro (RTT $2$ ms, $1$ Gbps) & $304$ ms & $80$ ms & $3.8\times$ \\
WAN (RTT $20$ ms, $1$ Gbps) & $409$ ms & $84$ ms & $4.9\times$ \\
WAN (RTT $80$ ms, $100$ Mbps) & $1156$ ms & $270$ ms & $4.3\times$ \\
\bottomrule
\end{tabular}
\end{table}

The garbled circuit wins on all three axes and at every network profile
(\Cref{tab:gc}): $9\times$ fewer rounds, $3.5\times$ less bandwidth, and
$3.8$--$4.9\times$ lower end-to-end latency --- not only on the WAN, as the
round-count argument would predict, but on the LAN as well. The enabler is
the honest-majority \emph{three-party} setting: a garbled circuit's input
labels are delivered by the co-garbler with no oblivious
transfer~\cite{ASTRA}, the expensive part of two-party garbling~\cite{Yao}
and absent here.

\Cref{tab:gc} is a bare comparison, which is not the pipeline's actual
nonlinearity. We therefore garble the \emph{entire} fused
ReLU$\circ$rescale, $\mathrm{out} = (v \geq 0)\,?\,(v \gg s):0$, as one
constant-round circuit, and run it as the real three-party protocol ---
two garblers, an evaluator, and a Y2A round returning the output to an
arithmetic sharing so the next scheme switch can consume it. Measured through
the same transport, that full protocol beats the \texttt{a2b} fused
nonlinearity by $1.7\times$ on the LAN and $1.4\times$ at RTT $80$ ms. We
flag the gap between this and the $2.7$--$3.4\times$ that the garbling
\emph{primitive} shows in isolation: the difference is the Y2A return leg and
the three-party label delivery, and it is the honest number for a protocol
one can deploy.

Integrated into the pipeline --- so that the fused ReLU, and through the same
mechanism every batched max-pool comparison, takes the Yao leg --- it takes
the MPC segment of the $28\times28$ CNN from $1.15$ s to $0.61$ s
($1.9\times$) and halves party-to-party traffic from $39.5$ to $21.8$ MB. The
gain survives real networks: on the trained medical model at RTT $20$ ms it
is $2.0\times$ on the segment ($3.39$ s to $1.70$ s), even though garbling
inverts the message \emph{count} ($88$ messages becomes $319{,}139$ small
ones while bytes halve). Both routes are bit-exact against the plaintext
reference and against each other, which is what licenses reading the ratio as
pure cost.

\paragraph{The last dependent reduction.} Batching a layer requires its
operations to be independent, and one gadget was not: \bt{argmax} folded
$\text{classes}-1$ comparisons left to right, each waiting on the previous
winner. It did less work than any other phase and, under WAN latency,
took the most time --- $1.31$ s of a $2.03$ s MPC segment at RTT $20$ ms.
Reassociating the fold into a \emph{tournament} --- pair adjacent candidates,
decide a whole level with one batched comparison, select value and index with
one batched multiplication --- makes it $\lceil\log_2 n\rceil$ batched levels
instead of $n-1$ sequential ones, for $2.4\times$ at RTT $20$ ms. The
transformation must preserve tie-breaking (the first maximum wins, as in the
plaintext reference), which it does because the bracket pairs adjacent
elements in index order; we verify it exhaustively over small alphabets
rather than on sampled inputs, since a tie split across sibling subtrees is
exactly the case random logits never produce.

\subsection{Convolution packing}

Replacing the im2col/Toeplitz convolution with the reversed-axis packing of
\Cref{sec:fhe} reduces the number of ciphertext multiplications from
hundreds to one per output channel, independent of image size. On the
$28\times28$ classifier the convolution's FHE time drops from $77.5$ s to
$0.26$ s ($294\times$), and the FHE share of runtime falls to a few percent
--- which is what moved the cost center to the MPC segment and motivated
\Cref{sec:mpc}.

That is no longer where it sits. With the segment's import batched and its
nonlinearities garbled, the same $28\times28$ classifier spends $1.35$ s of
its $1.99$ s in FHE and $0.61$ s in MPC, and the trained two-convolution
medical model of \Cref{sec:diagnosis} --- four encrypted linear layers to
this one's two --- spends $73\%$ of an online diagnosis in FHE. The cost center has
crossed back. We note this in favor of the construction rather than against
it: the switch-as-bootstrap argument buys unbounded depth at a
\emph{constant} FHE cost per layer, and the residual cost is now that
constant, paid four times, rather than any part of the switching machinery
the construction adds.

\subsection{Real MNIST: a fully encrypted two-convolution CNN}\label{sec:mnist}

We train the exact architecture the pipeline evaluates --- $\texttt{conv}(8,
3\times3)\to\ReLU\to\text{pool}\to\texttt{conv}(16,3\times3)\to\ReLU\to
\text{pool}\to\texttt{dense}(128)\to\ReLU\to\texttt{dense}(10)$ --- on the
full $60{,}000$-image MNIST~\cite{MNIST} training set, then \emph{post-train
quantize into the exact integer fixed-point scheme the encrypted pipeline
runs}. The reported accuracy is thus that of the deployed integer model, not
of a float model that would lose accuracy under quantization. This model has
four encrypted linear layers (two convolutions, two dense), three MPC
nonlinearities, and three scheme switches, and the second convolution is
multi-channel and packed with the same reversed-axis layout.

\begin{table}[t]
\centering
\caption{Full-test-set accuracy of the deployed integer model, by
architecture. All fit the depth-one noise budget ($\text{nbits}\le 30$). The
selected architecture (last row) was subsequently retrained for longer under
the same recipe, which is the $98.89\%$ reported; the other rows are from the
original sweep, so the column is a comparison of architectures, not of
training budgets.}
\label{tab:arch}
\begin{tabular}{@{}lr@{}}
\toprule
architecture & MNIST test accuracy \\
\midrule
$\texttt{conv}(8) \to \texttt{dense}(10)$ & $94.96\%$ \\
$\texttt{conv}(12) \to \texttt{dense}(10)$ & $95.73\%$ \\
$\texttt{conv}(8) \to \texttt{dense}(128) \to \texttt{dense}(10)$ & $97.63\%$ \\
$\texttt{conv}(8) \to \texttt{dense}(256)\to\texttt{dense}(128)\to\texttt{dense}(10)$ & $97.89\%$ \\
$\texttt{conv}(8)\to\texttt{conv}(16)\to\texttt{dense}(128)\to\texttt{dense}(10)$ & $\mathbf{98.89\%}$ \\
\bottomrule
\end{tabular}
\end{table}

Because the construction makes depth free, we used it as an instrument to
search architectures (\Cref{tab:arch}): a second convolution layer beats both
a wider single convolution and a deeper fully-connected stack --- spatial
depth is what MNIST rewards, and the construction supplies it at no
cryptographic cost. The chosen two-convolution model attains $98.89\%$ on the
$10{,}000$-image test set. The Rust integer forward pass agrees with an
independent numpy reference digit-for-digit, and the \emph{encrypted}
pipeline reproduces the plaintext prediction on every real digit tested, at
zero bootstraps. Because the hybrid is bit-exact, the encrypted accuracy
\emph{equals} the plaintext accuracy of $98.89\%$: privacy costs nothing in
accuracy, demonstrated on real inputs rather than asserted.

\subsection{Preprocessing the switch}\label{sec:preproc}

With the MPC segment down to a fraction of runtime, we broke the FHE side
down the same way. On the $28\times28$ classifier it divides as: encrypting
the model $218$ ms, the ct$\times$ct products $210$ ms, decrypting the masked
openings $217$ ms, re-encrypting MPC shares $50$ ms, encrypting the input
$17$ ms --- and encrypting the \emph{masks} $656$ ms, $48\%$ of the FHE total
and $39$ of the $52$ encryptions the inference performs.

Those masks are uniform random values. They depend on neither the image nor
the weights, only on the network's shape, so every one of them can be
generated and encrypted before the query exists. They must still be
\emph{fresh per query} --- two queries masked identically would reveal the
difference of their pre-activations --- so preprocessing does not remove the
work, it moves it off the critical path, which is precisely the deployment
question. A server idle between patients has the time; a patient waiting for
a diagnosis does not. Precomputing also lets the three parties' mask
ciphertexts be \emph{summed} in advance, since the evaluator only ever needs
their sum, so the online step becomes a single ciphertext addition in place
of three encryptions and three additions.

Measured on the $28\times28$ classifier, this takes the online time from
$1.98$ s to $1.31$ s ($1.51\times$) against $0.68$ s of offline work per
query --- the same total, differently placed. On the trained two-convolution
medical model the effect is larger, because it performs four scheme switches
rather than two: \Cref{tab:diagnosis}. The encrypted diagnosis is unchanged,
which we check by running the identical inference both ways and requiring the
same class.

\subsection{Private diagnosis: real medical images end to end}\label{sec:diagnosis}

Handwritten digits are a correctness and cost instrument, not an application.
The setting this construction is actually for --- both the input \emph{and}
the model encrypted --- is one where neither party may see the other's asset:
a hospital holding patient images and a vendor holding a proprietary
diagnostic model. We therefore run the identical pipeline, unchanged, on
MedMNIST v2~\cite{MedMNIST}, which distributes biomedical images at
$28\times28$: \texttt{PneumoniaMNIST} (pediatric chest X-ray, pneumonia
vs.\ normal) and \texttt{OrganAMNIST} ($11$-way organ classification from
abdominal CT). We use the two grayscale sets, since the convolution packing
assumes a single input channel.

\Cref{tab:diagnosis} reports the deployed integer model's accuracy and the
cost of one encrypted diagnosis. Because the hybrid is bit-exact, the
encrypted diagnosis \emph{equals} the plaintext one on every image tested, so
the accuracy column is simultaneously the encrypted accuracy: privacy costs
nothing in diagnostic quality. What it costs is $2.9$ s per chest X-ray
online, against $2.0$ s of query-independent preprocessing; of the online
time $73\%$ is FHE arithmetic and $0.76$ s is the entire MPC segment
(import $0.02$ s, ReLU $0.51$ s, max-pool $0.21$ s, argmax $0.01$ s), over
$27.4$ MB of party-to-party traffic and zero bootstraps. Placing the three
parties on a wide-area link (RTT $20$ ms, $1$ Gbps) costs $2.1$ s more, of
which the MPC segment is $0.9$ s: the pipeline is not round-bound, which is
what the constant-round nonlinearities buy.

For scale, this diagnosis cost $32.7$ s at the start of the optimization
sequence in \Cref{sec:eval} and $2.9$ s at the end, an $11\times$ reduction in
which the cryptographic construction never changed --- only the batching of
its round trips, the protocol chosen for its comparisons, and what was moved
off the critical path. Every intermediate version was bit-exact against the
same plaintext reference, which is what let the sequence be driven by
measurement rather than by argument.

\begin{table}[t]
\centering
\caption{Private diagnosis on MedMNIST v2. Accuracy is the deployed
\emph{integer} model on the full official test split; the encrypted pipeline
reproduces it bit-exactly, so it is also the encrypted accuracy. Latency is
one encrypted diagnosis with garbled nonlinearities, single-threaded, no GPU.}
\label{tab:diagnosis}
\begin{tabular}{@{}lrr@{}}
\toprule
 & \texttt{PneumoniaMNIST} & \texttt{OrganAMNIST} \\
\midrule
task & X-ray, 2 classes & CT, 11 classes \\
test images & $624$ & $17{,}778$ \\
majority-class baseline & $62.50\%$ & $18.48\%$ \\
accuracy (integer $=$ encrypted) & $\mathbf{84.78\%}$ & $\mathbf{81.18\%}$ \\
float parent of the same net & $84.78\%$ & $81.03\%$ \\
\midrule
latency, \texttt{a2b} nonlinearities & $5.71$ s & $5.95$ s \\
$+$ garbled nonlinearities & $4.94$ s & $5.05$ s \\
$+$ preprocessed masks (online) & $\mathbf{2.90}$ s & $\mathbf{3.00}$ s \\
\quad of which FHE & $2.13$ s & $2.13$ s \\
\quad of which MPC & $0.76$ s & $0.85$ s \\
offline, per query & $2.03$ s & $2.04$ s \\
\midrule
same online figure at WAN (RTT $20$ ms, $1$ Gbps) & $5.00$ s & $5.30$ s \\
\quad of which MPC & $1.71$ s & $2.02$ s \\
\midrule
FHE bootstraps & $0$ & $0$ \\
\bottomrule
\end{tabular}
\end{table}

Two things are worth reading off the table. First, \emph{quantization is
free}: the integer network the FHE layer requires lands within $0.2$ points
of its float parent, so the fixed-point scheme costs no accuracy and neither
does the encryption. Second, \emph{the model, not the cryptography, is the
accuracy ceiling}. A published ResNet-18 at $28\times28$ reaches $0.854$ and
$0.935$ on these two sets; ours is a $24$k-parameter two-convolution
classifier sized to the depth-one noise budget, and it is essentially at the
published number on pneumonia while giving up $12$ points to depth on organ
classification. Closing that gap is a capacity question, and the construction
makes depth cryptographically free --- but each additional layer costs one
more FHE multiplication, which \Cref{tab:diagnosis} shows is now the
dominant term.

\paragraph{Scope.} MedMNIST is downsampled and pre-processed; clinical
imaging is $512^3$ volumes, and nothing here is a clinical claim. The claim
is about \emph{cost structure} --- that a fully encrypted diagnosis, model
included, costs seconds rather than minutes, and that its cost is now
concentrated in a single primitive. Scaling past $28\times28$ is an argument
the packing constraint has to make analytically, not one this measurement
makes.

\subsection{Positioning against prior inference systems}

\Cref{tab:sota} places the result against representative two-party inference
systems. The residual $\sim\!1\%$ accuracy gap is model size, a training
question the construction is indifferent to. The distinguishing axis is the
last column: prior systems are client/server protocols in which the server
evaluates on \emph{plaintext} weights, so the model is not hidden; here every
linear layer is ciphertext$\times$ciphertext, so the weights are encrypted
too, and only the class index leaves the cryptographic domain.

\begin{table}[t]
\centering
\caption{Positioning on MNIST. The comparison axis is \emph{what is hidden},
not accuracy.}
\label{tab:sota}
\begin{tabular}{@{}lll@{}}
\toprule
system & MNIST acc & what is hidden \\
\midrule
\textbf{this work} & $\mathbf{98.0\%}$ & \textbf{image \emph{and} weights} (ct$\times$ct) \\
Gazelle~\cite{Gazelle} & ${\sim}99\%$ & image only (plaintext weights) \\
Delphi~\cite{Delphi} & ${\sim}99\%$ & image only \\
CrypTFlow2~\cite{CrypTFlow2} & ${\sim}99\%$ & image only \\
MiniONN~\cite{MiniONN}/Cheetah~\cite{Cheetah} & ${\sim}99\%$ & image only \\
\bottomrule
\end{tabular}
\end{table}

\section{Related Work}\label{sec:related}

\paragraph{Hybrid linear-FHE / nonlinear-MPC inference.}
Gazelle~\cite{Gazelle}, Delphi~\cite{Delphi}, Cheetah~\cite{Cheetah}, and
CrypTFlow2~\cite{CrypTFlow2} pioneered evaluating linear layers under
HE and nonlinearities under MPC. They treat the switch as engineering
overhead and keep the model in plaintext. We identify the switch as a
bootstrapping-elimination mechanism enabling unbounded depth with a
depth-one FHE engine, and we encrypt the weights.

\paragraph{Preprocessing.} Delphi~\cite{Delphi} moves query-independent HE
work off the critical path. We apply the same idea to the scheme switch
rather than to the linear layers: the masks are uniform and depend only on
the network's shape, so they are generated and encrypted before the query
arrives (\Cref{sec:preproc}). Where Delphi preprocesses to make the linear
layer cheap online, we preprocess the \emph{switching} machinery --- the part
a switch-as-bootstrap construction performs most often.

\paragraph{Scheme switching between FHE schemes.} CHIMERA~\cite{CHIMERA} and
PEGASUS~\cite{PEGASUS} switch FHE$\leftrightarrow$FHE to pick the cheapest
scheme per operation. We switch FHE$\leftrightarrow$MPC and exploit the
re-encryption as a noise refresh.

\paragraph{Leveled and bootstrapped FHE.} BGV~\cite{BGV}, BFV~\cite{FV}, and
CKKS~\cite{CKKS} evaluate bounded-depth circuits and bootstrap (following
Gentry~\cite{Gentry}) to go deeper; TFHE~\cite{TFHE} bootstraps every gate.
Our construction removes the bootstrap for the deep-inference workload by
interleaving an MPC round trip that already exists for the nonlinearity.

\paragraph{Overstretched NTRU.} The insecurity of NTRU at large modulus is
established by~\cite{ABDucas,KirchnerFouque} and sharpened
by~\cite{DucasvanWoerden}; we rely on it to justify using NTRU only as a
non-secure reference and RLWE as the secure instantiation.

\section{Limitations}\label{sec:limitations}

The construction wins on depth scaling, on encrypting the model, and on the
thin-pipe regime, \emph{not} on a single small layer over a fast LAN, where
pure MPC is faster. Absolute latency remains seconds per image and is now
dominated by the FHE arithmetic --- $73\%$ of an online diagnosis, and the
part of it that preprocessing cannot touch, since the products depend on both
the image and the weights. Our FHE engine is
a single-threaded schoolbook implementation with no GPU and no
machine-specific NTT; the residual cost is therefore an engineering number,
not an information-theoretic one, but it is the number that matters for
deployment and we do not claim to have optimized it.
Softmax and layer normalization (as in transformers) need additional MPC
subprotocols (exponentiation, reciprocal, inverse square root) that we do not
implement; deep MLP/CNN ReLU stacks are supported. Layer width per ciphertext
is bounded by the ring dimension through the packing constraint; wider layers
chunk. Multi-channel inputs beyond the ring-dimension bound (e.g. CIFAR-scale
three-channel images) require splitting channels across ciphertexts or a
larger ring --- a packing question, not a limit of the construction.

\section{Conclusion}\label{sec:conclusion}

The scheme switch that hybrid private-inference systems already perform for
each nonlinearity is, at the same time, a noise reset: it re-encrypts the
activation to a fresh, depth-one ciphertext. Reading the switch this way
turns bootstrapping from a per-depth cost into an absent one, and lets the
FHE engine be a minimal, evaluation-key-free, single-multiply primitive. We
gave the argument precisely, instantiated it with both operands encrypted,
built a triple-free batched MPC layer with constant-round garbled
nonlinearities, and validated the whole pipeline bit-exactly --- culminating
in a fully encrypted two-convolution CNN that diagnoses a real pediatric
chest X-ray in $4.9$ seconds, model included, with zero FHE bootstraps and a
diagnosis provably identical to the plaintext one.

Getting there also cost us a wrong answer worth reporting. Having identified
the MPC segment as the bottleneck, we spent our optimization effort on MPC
protocols; the segment's cost turned out to be an unbatched round trip in the
scheme-switch import --- the seam between the two representations, which a
construction organized around switching crosses more often than anything
else. Removing it was an $8.3\times$ end-to-end win and moved the cost center
back onto the FHE multiplication, which is where the depth-one design says it
should be. We suspect the seam is systematically under-measured in hybrid
systems, because it belongs to neither of the two literatures whose
primitives it joins.

\bibliographystyle{plain}
\bibliography{references}

\end{document}
